\section*{अध्याय \ref{ch:b2:cell-differentiation} --- कोशिका विभेदन: कंकाल पेशी कोशिका}

\begin{solution}{exo:b2:cell-differentiation:1}
\omterm{def:b2:cell-differentiation:differentiation}{विभेदन}: किसी एक कोशिका-प्रकार के जीनों तथा कार्यों की स्थिर
अभिव्यक्ति। निर्धारण: किसी नियति के अभिव्यक्त होने से पहले उसके प्रति
प्रतिबद्धता। \omterm{def:b2:cell-differentiation:differentiation}{शक्तता}: अब भी खुली नियतियों की परास। युग्मनज:
पूर्णशक्त; \omterm{def:b2:mammal-reproduction:implantation}{अंतःकोशिका पिंड} की कोशिका: बहुशक्त; उपग्रही कोशिका:
एकशक्त (केवल पेशी) --- फिर भी एक \omterm{def:b2:cell-differentiation:differentiation}{स्तंभ कोशिका}; \omterm{def:b2:cell-differentiation:myogenesis}{पेशी तंतु}:
विभेदित, विभाजनोत्तर, कोई \omterm{def:b2:cell-differentiation:differentiation}{शक्तता} नहीं।
\end{solution}

\begin{solution}{exo:b2:cell-differentiation:2}
त्वक्-पेशीखंड की कोशिका (Pax3) $\to$ प्रेरित \omterm{def:b2:cell-differentiation:myogenesis}{पेशीकोरक} (\omterm{def:b2:cell-differentiation:myogenesis}{MyoD}, Myf5),
जो FGF रहते विभाजित होता है $\to$ चक्र से निकास, मायोजेनिन $\to$ पंक्तिबंधन
और संलयन से \omterm{def:b2:cell-differentiation:myogenesis}{पेशीनलिका} $\to$ पेशी-जीन चालू, \omterm{def:b2:cell-differentiation:fibre}{पेशीतंतुक}
जुड़े (MRF4) $\to$ परिधीय केंद्रकों वाला परिपक्व तंतु, और साथ में
निष्क्रिय उपग्रही कोशिकाएँ (Pax7)।
\end{solution}

\begin{solution}{exo:b2:cell-differentiation:3}
यह कि किसी विभेदित कोशिका का केंद्रक पूरा प्राणी बनाने को आवश्यक हर जीन
रखता है, और यह कि अंडाणु का कोशिकाद्रव्य उसका कार्यक्रम फिर से बिठा
सकता है: यानी \omterm{def:b2:cell-differentiation:differentiation}{विभेदन} प्रतिवर्ती है और DNA का लोप नहीं। उसने यह
नहीं दिखाया कि हर विभेदित केंद्रक फिर बिठाया जा सकता है (अधिकांश
प्रत्यारोपण विफल हुए), न यह कि वह पुनर्बिठान कैसे चलता है, न यह कि वह
विभेदित कोशिका स्वयं अपनी जगह नियति बदल सकती है।
\end{solution}

\begin{solution}{exo:b2:cell-differentiation:4}
\omterm{def:b2:cell-differentiation:fibre}{सार्कोमियर} को घेरती Z चक्रिकाएँ; उन पर टिके पतले ऐक्टिन तंतु;
M रेखा पर केंद्रित मोटे मायोसिन तंतु; A पट्टी
(मोटे तंतु) और I पट्टी (केवल पतले)। छोटा होने पर पतले
तंतु M रेखा की ओर सरकते हैं और Z चक्रिकाएँ एक-दूसरे के पास आ जाती हैं:
I पट्टी और H क्षेत्र सँकरे होते हैं, A पट्टी --- यानी मोटे तंतु की
लंबाई --- नहीं बदलती, और कोई भी तंतु छोटा नहीं होता।
\end{solution}

\begin{solution}{exo:b2:cell-differentiation:5}
$0.2/2.5\times 10^{-6} = \num{80000}$ \omterm{def:b2:cell-differentiation:fibre}{सार्कोमियर}; $0.2/25\times 10^{-6}
= \num{8000}$ केंद्रक; प्रति केंद्रक एक \omterm{def:b2:cell-differentiation:myogenesis}{पेशीकोरक}, इसलिए लगभग \num{8000}
\omterm{def:b2:cell-differentiation:myogenesis}{पेशीकोरक} संलयित हुए (उपग्रही कोशिकाएँ बाद में और जोड़ती हैं)।
\end{solution}

\begin{solution}{exo:b2:cell-differentiation:6}
FGF Id प्रेरित करती है, जो वह साथी प्रोटीन बाँध लेती है जिसे \omterm{def:b2:cell-differentiation:myogenesis}{MyoD} को कोई
DNA-बंधक द्विलक बनाने के लिए चाहिए; \omterm{def:b2:cell-differentiation:myogenesis}{MyoD} उपस्थित तो है पर निठल्ला, और वे
कोशिकाएँ विभाजित होती रहती हैं। FGF हटते ही Id क्षीण होती है, \omterm{def:b2:cell-differentiation:myogenesis}{MyoD}
द्विलक बनाता है, मायोजेनिन चालू हो जाती है, और वे कोशिकाएँ चक्र से निकलकर
संलयित हो जाती हैं। जो \omterm{def:b2:cell-differentiation:myogenesis}{पेशीकोरक} Id लगातार बनाते हों वे कभी संलयित नहीं
होते: वे अनिश्चित काल तक विभाजित होते रहते हैं और कोई पेशी नहीं बनाते।
\end{solution}

\begin{solution}{exo:b2:cell-differentiation:7}
शर्त $0.4m = m^{2}/(1 + m^{2}) + 0.02$। $m = 2.1$ पर: बायाँ
$0.84$, दायाँ $4.41/5.41 + 0.02 = 0.835$: यानी एक स्थायी अवस्था। नीची
अवस्था: छोटे $m$ के लिए $0.4m \approx m^{2} + 0.02$, इसलिए $m \approx
0.055$ (बायाँ $0.022$, दायाँ $0.023$)। 1 तक चढ़ाने पर, जो
$0.41$ के पास की देहली से ऊपर है, वह कोशिका $2.1$ तक चढ़ जाती है और वहीं
रहती है: यानी वह प्रतिबद्ध है।
\end{solution}

\begin{solution}{exo:b2:cell-differentiation:8}
$k = 0.6$ के साथ: $m = 1$ पर विघटन $0.6$ संश्लेषण
$0.52$ से अधिक है; $m = 0.5$ पर $0.30 > 0.22$; और $m = 2$ पर $1.2 > 0.82$: यानी
नीची अवस्था से आगे हर जगह वह रेखा उस वक्र के ऊपर पड़ी है, और वही
एकमात्र स्थायी अवस्था है। जिस कोशिका की \omterm{def:b2:cell-differentiation:myogenesis}{MyoD} बहुत तेज़ विघटित होती है
वह चालू अवस्था थाम ही नहीं सकती: वह प्रेरण चाहे जैसा भी हो, अविभेदित
अवस्था में लौट जाती है --- यानी कोई पेशी नहीं।
\end{solution}

\begin{solution}{exo:b2:cell-differentiation:9}
(क) कोई \omterm{def:b2:cell-differentiation:myogenesis}{पेशीकोरक} नहीं, कोई कंकाल पेशी नहीं: वे दोनों कारक
निर्धारण के लिए अतिरेकी हैं, इसलिए दोनों का जाना ज़रूरी है। (ख) \omterm{def:b2:cell-differentiation:myogenesis}{पेशीकोरक}
बनते और पंक्ति में लगते हैं पर कभी संलयित नहीं होते और न पेशी-प्रोटीन
बनाते हैं: \omterm{def:b2:cell-differentiation:differentiation}{विभेदन} के लिए मायोजेनिन आवश्यक है। (ग) लगभग सामान्य, क्योंकि
निर्धारण में Myf5 \omterm{def:b2:cell-differentiation:myogenesis}{MyoD} की जगह ले लेती है।
\end{solution}

\begin{solution}{exo:b2:cell-differentiation:10}
किसी तंतुकोरक में पेशी-जीन ऐसे क्रोमैटिन में पड़े हैं जिसे \omterm{def:b2:cell-differentiation:myogenesis}{MyoD} खोल
सकता है; जबकि किसी यकृत कोशिका में वे विषमक्रोमैटिन में बँधे हो सकते
हैं, या यकृत के अपने \omterm{prop:b2:cell-differentiation:transcriptional}{प्रमुख नियामक} उन्हें दबा सकते हैं, इसलिए वह कोशिका
सक्षम नहीं है। परीक्षण: \omterm{def:b2:cell-differentiation:myogenesis}{MyoD} डालने से पहले यकृत कोशिकाओं को क्रोमैटिन
ढीला करने वाली किसी औषधि (कोई हिस्टोन डीऐसीटिलेज़ संदमक) से उपचारित
कीजिए, या \omterm{def:b2:cell-differentiation:myogenesis}{MyoD} के साथ ऐसा कोई कारक डालिए जो यकृत का कार्यक्रम हटा दे,
और मायोसिन की अभिव्यक्ति गिनिए।
\end{solution}

\begin{solution}{exo:b2:cell-differentiation:11}
लंबे समय तक कम बारंबारता का स्फुरण अंतःकोशिकीय कैल्सियम को
देर तक चढ़ाए रखता है; और कैल्सियम-सक्रियित पथ मंद मायोसिन,
सूत्रकणिका-जनन, मायोग्लोबिन तथा केशिका-वृद्धि के जीन चालू कर देते हैं
और तीव्र समरूप बंद। उस तंतु का प्रमुख
कार्यक्रम (\omterm{def:b2:cell-differentiation:myogenesis}{MyoD} अवस्था, \omterm{def:b2:cell-differentiation:fibre}{सार्कोमियर} की स्थापत्य-रचना) अछूता रहता है:
प्रशिक्षण यह बदलता है कि पेशी-जीनों में से कौन अभिव्यक्त हों, उस कोशिका
की पहचान नहीं, और किसी तंत्रिका का कोई संकेत वह उपास्थि-कार्यक्रम फिर
नहीं खोल सकता जो निर्धारण पर बंद हो गया था।
\end{solution}

\begin{solution}{exo:b2:cell-differentiation:12}
किसी स्वयं-सक्रियकारी कारक की चालू अवस्था किसी प्रतिपुष्टि चक्र का
स्थायी बिंदु है, जो लगातार संश्लेषण से बना रहता है, जीनों के किसी
परिवर्तन से नहीं; और क्रोमैटिन के चिह्न एक दूसरी परत जोड़ देते हैं जो
चुप किए गए जीन बंद रखती है। इसलिए पुनर्बिठान के लिए उस परिपथ को उसकी
देहली से नीचे धकेलना और साथ ही वह क्रोमैटिन फिर खोलना पड़ता है, जो
अंडाणु का कोशिकाद्रव्य या चार कारक कोशिकाओं के किसी छोटे अंश में ही कर
पाते हैं: संभव, क्योंकि कुछ खोया नहीं है, पर अदक्ष, क्योंकि
दो स्वतंत्र ताले एक साथ खोलने पड़ते हैं।
\end{solution}

\begin{solution}{pb:b2:cell-differentiation:1}
\textbf{1.} $0.2/2.5\times 10^{-6} = \num{80000}$ \omterm{def:b2:cell-differentiation:fibre}{सार्कोमियर};
$0.2/25\times 10^{-6} = \num{8000}$ केंद्रक।
\textbf{2.} प्रति तंतु लगभग \num{8000}; और पूरी पेशी के लिए $4\times 10^{9}$।
\textbf{3.} प्रति तंतु $\pi(30\times 10^{-6})^{2}\times 0.2 = \qty{5.7e-10}{m^3}$;
और पूरी पेशी के लिए $2.8\times 10^{-4}$ \unit{m^3}, यानी \qty{0.28}{L}।
\textbf{4.} $5.7\times 10^{-10}/8000 = \qty{7e-14}{m^3} =
\qty{70000}{\micro m^3}$: यानी प्रति केंद्रक 35 साधारण कोशिकाओं जितना।
\textbf{5.} भीतरी भाग \omterm{def:b2:cell-differentiation:fibre}{पेशीतंतुकों} से ठसाठस भरा है जिन्हें सिरे से सिरे
तक बिना टूटे चलना पड़ता है; और किनारे के केंद्रक उस जालक को
अटूट छोड़ देते हैं तथा झिल्ली और उसके संकेतों के पास बैठते हैं।
\textbf{6.} $4\times 10^{9}/2000 = 2\times 10^{6}$: यानी 21 दुगुनेपन, 10.5
दिन।
\textbf{7.} $0.4m = m^{2}/(1 + m^{2}) + 0.02$। $m = 0.055$ पर: बायाँ
$0.022$, दायाँ $0.003 + 0.02 = 0.023$।
\textbf{8.} $m = 2.1$ पर: बायाँ $0.84$, दायाँ $0.815 + 0.02 = 0.835$।
\textbf{9.} $m = 0.41$ पर: बायाँ $0.164$, दायाँ $0.144 + 0.02 = 0.164$।
उससे ठीक नीचे संश्लेषण विघटन से कम है, इसलिए $m$ गिरकर
$0.055$ की ओर जाता है; और ठीक ऊपर संश्लेषण विघटन से अधिक है और $m$
$2.1$ की ओर चढ़ता है: यानी कोई भी विस्थापन बढ़ता जाता है --- अस्थिर।
\textbf{10.} $0.6 > 0.41$: पहली कोशिका चालू अवस्था पर चली जाती है और
विभेदित हो जाती है; $0.3 < 0.41$: दूसरी ढलकर $0.055$ पर आ जाती है और
नहीं होती।
\textbf{11.} $b = 0.4$ के साथ संश्लेषण-वक्र $0.4$ से आरंभ होता है और
रेखा $0.4m$ के ऊपर $m \approx 3.3$ तक बना रहता है: यानी नीचा कटान
जाता रहा। वह कोशिका बिना किसी प्रेरण के स्वतःस्फूर्त विभेदित हो जाती है।
\textbf{12.} विभाजन $m$ आधा कर देता है; पुत्रियों को $2.1/2 =
1.05$ विरासत में मिलता है, जो देहली $0.41$ से ऊपर है, और वे फिर $2.1$ तक
चढ़ जाती हैं: यानी वह अवस्था हर पुत्री में फिर स्थापित हो जाती है।
प्रतिबद्धता तब तक बची रहती है जब तक आधा किया गया स्तर उस देहली से ऊपर रहे।
\textbf{13.} प्रति मिमी $4$ $\times$ \qty{200}{mm} $= 800$ उपग्रही
कोशिकाएँ प्रति तंतु; और \num{8000} केंद्रकों का \qty{10}{\%}, यानी
\num{800}, बदलना पड़ता है।
\textbf{14.} घायल खंड में $80$ उपग्रही कोशिकाएँ हैं; $800/80 =
10$, इसलिए 4 दुगुनेपन ($\times 16$, यानी \num{1280}: 800 संलयित होती हैं और 480
भंडार भरने को बचती हैं) --- यानी \qty{72}{h}, तीन दिन।
\textbf{15.} विभाजन के तीन दिन बनाम प्रेक्षित दो से तीन सप्ताह:
बाक़ी में भक्षकाणुओं द्वारा मलबे की सफ़ाई, सक्रियण का विलंब, संलयन,
\omterm{def:b2:cell-differentiation:fibre}{पेशीतंतुकों} का जुड़ना, पुनर्तंत्रिकन और नए खंड का
परिपक्वन सम्मिलित है।
\textbf{16.} सक्रियित कोशिकाओं में \omterm{def:b2:cell-differentiation:myogenesis}{MyoD} (और मायोजेनिन) को पेशी-जीन फिर
चालू करने ही पड़ते हैं; और विभाजन तथा \omterm{def:b2:cell-differentiation:differentiation}{विभेदन} परस्पर अपवर्जी हैं
--- मायोजेनिन के काम करने और संलयन होने से पहले Id को गिरना और चक्र को
रुकना पड़ता है।
\textbf{17.} $0.95^{n} = 0.5$: $n = 13.5$, यानी लगभग चौदह चोटें।
\textbf{18.} $800\times 0.95^{20} = 800\times 0.36 \approx 290$।
\textbf{19.} हर संकुचन तंतु फाड़ता है, हर फटन उस भंडार को पुकारती है,
और वह भंडार हर बार कुछ प्रतिशत सिकुड़ता है; वर्षों बाद वह चुक जाता है,
फटे तंतु फिर नहीं गढ़े जाते, और तंतुकोरक उन ख़ाली जगहों को कोलैजन से भर देते हैं।
\textbf{20.} $(80/60)^{2} = 1.78$; केंद्रक $8000\times 1.78 = \num{14200}$:
यानी उपग्रही कोशिकाओं से लगभग \num{6200} अतिरिक्त।
\textbf{21.} $0.28\times 1.78 = \qty{0.5}{L}$।
\textbf{22.} बदला: मायोसिन समरूपों के जीन, सूत्रकणिका तथा
केशिका घनत्व, मायोग्लोबिन। नहीं बदला: पेशी-पहचान (\omterm{def:b2:cell-differentiation:myogenesis}{MyoD}
अवस्था), \omterm{def:b2:cell-differentiation:fibre}{सार्कोमियर} की स्थापत्य-रचना, केंद्रक और उनका जीनोम।
\textbf{23.} तंतु विभाजनोत्तर हैं और गुणित नहीं हो सकते; एकमात्र
रास्ते मोटे तंतु और उपग्रही कोशिकाओं से प्रति तंतु अधिक केंद्रक हैं।
\textbf{24.} प्रशिक्षण: सीमित मोटापन, क्योंकि वह तंतु केंद्रक नहीं जोड़
सकता; चोट: कोई मरम्मत नहीं --- नष्ट हुए खंडों का स्थान
तंतुमय ऊतक ले लेता है, जैसा दुष्पोषण में।
\textbf{25.} प्रति तंतु \num{8000} केंद्रक; देहली $m \approx 0.41$;
और किसी मरम्मत के लिए उपग्रही-कोशिका विभाजन के लगभग तीन दिन।
\end{solution}
